\documentclass[10pt,UTF8]{article}

\usepackage{geometry}

\usepackage{ctex}

\geometry{a4paper,scale=0.9}
\ctexset{today=old}

\usepackage[bookmarksnumbered=true]{hyperref}  % 生成大纲


\usepackage{times}

% \usepackage{makecell}
% \usepackage{multirow}

% \usepackage{graphicx} %插入图片的宏包
% \usepackage{subfigure}
% \usepackage{float} %设置图片浮动位置的宏包

\usepackage{amsmath}
\usepackage{amssymb}
\usepackage{mathtools}

\usepackage{physics}

% \usepackage{siunitx}
% % 关闭警告
% \ExplSyntaxOn
% \msg_redirect_name:nnn { siunitx } { physics-pkg } { none }
% \ExplSyntaxOff

%%%%%%%%%%%%%%%%%%%%%%%% tcolorbox %%%%%%%%%%%%%%%%%%%%%%%%%%
\usepackage[most]{tcolorbox}

\newenvironment{tipbox}
{\begin{tcolorbox}[colback=green!5!white,colframe=green!75!black]}
{\end{tcolorbox}}

\newenvironment{conceptbox}
{\begin{tcolorbox}[colback=blue!5!white,colframe=blue!75!black]}
{\end{tcolorbox}}

\newenvironment{thmbox}
{\begin{tcolorbox}[colback=yellow!5!white,colframe=yellow!75!black]}
{\end{tcolorbox}}

\newenvironment{definition box}
{\begin{tcolorbox}[colback=blue!5!white,colframe=blue!75!black,parbox=false,before upper=\par,before lower=\par]}
{\end{tcolorbox}}

\newenvironment{theorem box}
{\begin{tcolorbox}[colback=yellow!5!white,colframe=yellow!75!black,parbox=false,before upper=\par,before lower=\par]}
{\end{tcolorbox}}
%%%%%%%%%%%%%%%%%%%%%%%%%%%%%%%%%%%%%%%%%%%%%%%%%%%%%%%%%%%%
            
%%%%%%%%%%%%%%%%%%%%% amsthm %%%%%%%%%%%%%%%%
\usepackage{amsthm}

\theoremstyle{definition}
\newtheorem{definition inner}{Definition}[section]
\newenvironment{definition}[1][]
{\begin{definition box}\begin{definition inner}[#1]\hfill\par}
{\end{definition inner}\end{definition box}}

\theoremstyle{plain}
\newtheorem{theorem inner}{Theorem}[section]
\newenvironment{theorem}[1][]
{\begin{theorem box}\begin{theorem inner}[#1]\hfill\par}
{\end{theorem inner}\end{theorem box}}

\theoremstyle{definition}
\newtheorem*{proof inner}{\textit{Proof}}
\renewenvironment{proof}[1][]
{\begin{proof inner}[#1]\hfill\par}
{\qed\end{proof inner}}

\theoremstyle{remark}
\newtheorem*{remark}{\textit{Remark}}
%%%%%%%%%%%%%%%%%%%%%%%%%%%%%%%%%%%%%%%%%%%%%%



\newcommand{\mycases}[2][0.8]{
    \left\{\vphantom{\scalebox{#1}{
        $\begin{matrix}#2\end{matrix}$
    }}\right.\!\!\!\!
    \begin{array}{l@{\quad}l}#2\end{array}
}

\newcommand{\e}{\mathrm{e}{\mkern 1 mu}}


\newcommand{\perm}[2]{\ensuremath{{}_{#1}P_{#2}}}  % 排列数
\newcommand{\comb}[2]{\ensuremath{{}_{#1}C_{#2}}}  % 组合数
\newcommand{\Prob}{\mathbb{P}\qty}


\begin{document}

\part*{Conditional Probability and Bayes' Theorem}

\section{Conditional Probability}

\subsection{Definition}

\begin{definition}[Conditional Probability]
    The conditional probability of an event $A$, given that
    the event $B$ has occurred, is defined by
    \begin{equation*}
        \Prob(A\mid B) := \frac{\Prob(A\cap B)}{\Prob(B)}
    \end{equation*}
    provided that $\Prob(B)\ne 0$.
\end{definition}

\begin{remark}
    It can be easily verified that the \emph{conditional probability} defined above
    satisfies the \emph{probability axioms}.
    Thus it is actually a \emph{probability function} whose \emph{sample space} is $B$.
    
    The essence of conditional probability is that
    it is a corresponding probability function
    of a given probability function.
\end{remark}

\subsection{Multiplication Rule}

\begin{definition}[Multiplication Rule]
    The probability that two events, $A$ and $B$, both occur is given
    by the multiplication rule
    \begin{gather*}
        \Prob(A\cap B) = \Prob(A)\Prob(B\mid A) \qq{provided} \Prob(A) > 0 \\
        \shortintertext{or}
        \Prob(A\cap B) = \Prob(B)\Prob(A\mid B) \qq{provided} \Prob(B) > 0
    \end{gather*}
\end{definition}


\section{Independent Events}

\subsection{Definition}
\begin{definition}[Independent Events]
    Events $A$ and $B$ are \textbf{independent} iff
    \begin{equation*}
        \Prob(A\cap B) = \Prob(A)\Prob(B)
    \end{equation*}
    Otherwise, they are called \textbf{dependent} events.
\end{definition}
    
\subsection{Consequence}

\paragraph{I}
When $\Prob(A)>0$, we have
\begin{equation*}
    \Prob(B\mid A) = \Prob(B) \quad \iff \quad A \text{ and } B \text{ are independent}
\end{equation*}

\paragraph{II}
When $\Prob(B)>0$, we have
\begin{equation*}
    \Prob(A\mid B) = \Prob(A) \quad \iff \quad A \text{ and } B \text{ are independent}
\end{equation*}

\paragraph{III}
$A$ and $B$ are independent, iff the following pairs of events are also independent:
\begin{itemize}
    \item $A$ and $B'$
    \item $A'$ and $B$
    \item $A'$ and $B'$
\end{itemize}

\subsection{Mutually Independent Events}

\begin{definition}[Three Mutually Independent Events]
    Events $A$, $B$, and $C$ are \textbf{(mutually) independent}
    iff the following two conditions hold:
    \begin{itemize}
        \item $A$, $B$, and $C$ are \textbf{pairwise independent}, i.e.,
        $A$ and $B$, $B$ and $C$, and $A$ and $C$ are independent, respectively.
        \item $\Prob(A\cap B\cap C) = \Prob(A)\Prob(B)\Prob(C)$
    \end{itemize}
\end{definition}

\begin{definition}[Mutually Independence for More Events]
    $A_1, A_2, \cdots, A_n$ are \textbf{(mutually) independent} iff
    \begin{equation*}
        \Prob(\bigcap_{i=1}^k A_i) = \prod_{i=1}^k \Prob(A_i)
    \end{equation*}
    holds for all $k=2,3,\cdots,n$.
\end{definition}

\section{Bayes' Theorem}

\begin{theorem}[Bayes' Theorem]
    \paragraph*{}Assume that
    \begin{enumerate}
        \item $B_1, B_2, \cdots, B_n$ are \textbf{mutually exclusive} and \textbf{exhaustively} events.
        \item The \textbf{prior probability} of all $B_i$ is nonzero, i.e.,
        \begin{equation*}
            \Prob(B_i) > 0, \quad i = 1, 2, \cdots, n
        \end{equation*}
    \end{enumerate}

    \paragraph*{}Then, for any event $A$, we have
    \begin{itemize}
        \item \textbf{Law of Total Probability}:
        \begin{equation*}
            \Prob(A) = \sum_{i=1}^n \Prob(A \cap B_i) = \sum_{i=1}^n \Prob(A\mid B_i) \Prob(B_i)
        \end{equation*}
        \item \textbf{Bayes' Rule}: If $\Prob(A)>0$, then we have
        the \textbf{posterior probability} of $B_k$:
        \begin{equation*}
            \Prob(B_k\mid A) = \frac{\Prob(B_k\cap A)}{\Prob(A)} = 
            \frac{\Prob(A\mid B_k) \Prob(B_k)}{\displaystyle\sum_{i=1}^n \Prob(A\mid B_i) \Prob(B_i)}
        \end{equation*}
    \end{itemize}
\end{theorem}

\end{document}
